\chapter{Aftrekken}\label{ch:g2:subtraction}

Als er bovenaan te weinig eenheden staan om van af te trekken, breken
we een tiental open --- het \emph{lenen}, de spiegel van het
onthouden. Dit hoofdstuk legt het uit met stokjes en oefent daarna het
cijferend aftrekken, met doortellen als handige kortere weg.

\section{Waarom we lenen}

\begin{example}[Een tiental openbreken]\label{ex:g2:subtraction:unbundle}
$42 - 17$: we willen $7$ eenheden weghalen bij $2$ eenheden --- te
weinig! Maak dus van één van de $4$ tientallen tien losse eenheden:
bovenaan staan nu $3$ tientallen en $12$ eenheden. Wegnemen:
$12 - 7 = 5$ eenheden en $3 - 1 = 2$ tientallen. Dus $42 - 17 = 25$.
\end{example}

\begin{omfigure}
\begin{tikzpicture}[scale=0.75]
  % een gebundeld tiental, links
  \foreach \i in {0,...,9} \draw[omExo, very thick]
    (0.2*\i,0) -- (0.2*\i,1.7);
  \draw[omExo, thick, rounded corners=3pt]
    (-0.2,0.62) rectangle (2.0,1.05);
  \node[below, font=\small] at (0.9,-0.25) {één tiental};
  % en rechts opengebroken in tien losse eenheden
  \draw[-Stealth, omThm, very thick] (2.6,0.85) -- (3.6,0.85);
  \foreach \i in {0,...,9} \draw[omDef, very thick]
    (4.2+0.34*\i,0) -- (4.2+0.34*\i,1.7);
  \node[below, font=\small] at (5.73,-0.25) {tien losse eenheden};
  \node[right, font=\small, align=left] at (8.0,0.85)
    {breek een tiental open om genoeg\\eenheden te hebben om weg te nemen};
\end{tikzpicture}

{\small Lenen is openbreken: één tiental wordt tien eenheden, zodat de
aftrekking door kan gaan.}
\end{omfigure}

\section{Het cijferen}

\begin{method}[Onder elkaar aftrekken]\label{met:g2:subtraction:column}
\begin{enumerate}
  \item Schrijf het grootste getal bovenaan, rechts uitgelijnd;
  \item trek de eenheden af; is het bovenste cijfer te klein, leen dan
        een tiental (dat wordt tien eenheden in deze kolom);
  \item trek de tientallen af, dan de honderdtallen;
  \item \emph{controleer}: de uitkomst plus het getal dat je hebt
        weggenomen moet het bovenste getal teruggeven.
\end{enumerate}
\end{method}

\begin{example}\label{ex:g2:subtraction:worked}
Reken $503 - 168$ uit:
\[
\begin{array}{r}
5\,0\,3 \\
-\ 1\,6\,8 \\
\hline
3\,3\,5
\end{array}
\]
Eenheden: $3 - 8$ is te klein; leen via de tientallen (die zijn er
niet, dus leen eerst bij de honderdtallen): bovenaan staan nu $4$
honderdtallen, $9$ tientallen en $13$ eenheden. Dan $13 - 8 = 5$;
$9 - 6 = 3$; $4 - 1 = 3$. Controle: $335 + 168 = 503$. \checkmark
\end{example}

\section{Doortellen}

\begin{method}[Aftrekken door door te tellen]\label{met:g2:subtraction:countup}
Liggen de getallen dicht bij elkaar, klim dan van het kleine naar het
grote en tel de sprongen op:
\[
83 - 78: \quad 78 \to 80 \text{ is } 2, \ 80 \to 83 \text{ is } 3,
\text{ dus } 83 - 78 = 5 .
\]
Doortellen is sneller dan lenen als de twee getallen dicht bij elkaar
liggen --- en het is precies hoe een winkelier wisselgeld teruggeeft.
\end{method}

\begin{omfigure}
\begin{tikzpicture}[scale=0.85]
  \draw[-Stealth] (76.6,0) -- (84.4,0);
  \foreach \i in {77,...,84} \draw (\i,-0.1) -- (\i,0.1);
  \foreach \i in {78,80,83}
    \node[below, font=\footnotesize] at (\i,-0.12) {\i};
  \draw[-Stealth, omDef, thick] (78,0.3) .. controls (79,0.9) ..
    (80,0.3);
  \node[omDef, font=\small] at (79,1.0) {$+2$};
  \draw[-Stealth, omThm, thick] (80,0.3) .. controls (81.5,1.05) ..
    (83,0.3);
  \node[omThm, font=\small] at (81.5,1.1) {$+3$};
\end{tikzpicture}

{\small $83 - 78$ met doortellen: twee sprongen, $2 + 3 = 5$.}
\end{omfigure}

\begin{example}[Wisselgeld geven]\label{ex:g2:subtraction:change}
Een stuk speelgoed kost $34$; je betaalt met $50$. Wisselgeld: tel door,
$34 \to 40$ is $6$, $40 \to 50$ is $10$, dus je krijgt $16$ terug. Dat
is de aftrekking $50 - 34 = 16$, op de manier van de winkelier.
\end{example}

\section{Oefeningen}

\begin{exercise}[$\star$]\label{exo:g2:subtraction:1}
Reken uit (zonder lenen): $58 - 23$; \quad $176 - 45$; \quad
$389 - 164$.
\end{exercise}

\begin{exercise}[$\star$]\label{exo:g2:subtraction:2}
Reken uit met lenen: $52 - 17$; \quad $73 - 28$; \quad $60 - 24$.
\end{exercise}

\begin{exercise}[$\star$]\label{exo:g2:subtraction:3}
Trek onder elkaar af en controleer elke uitkomst: $324 - 167$; \quad
$500 - 236$; \quad $703 - 158$.
\end{exercise}

\begin{exercise}[$\star$]\label{exo:g2:subtraction:4}
Reken uit door door te tellen: $62 - 58$; \quad $91 - 85$; \quad
$100 - 96$.
\end{exercise}

\begin{exercise}[$\star$]\label{exo:g2:subtraction:5}
Reken het wisselgeld uit (tel door): je betaalt $50$ voor een boek van
$42$; $20$ voor speelgoed van $13$; $100$ voor een tas van $75$.
\end{exercise}

\begin{exercise}[$\star$]\label{exo:g2:subtraction:6}
``Een boom is $156$ cm hoog. Dit jaar is hij $28$ cm gegroeid. Hoe hoog
was hij vorig jaar?'' Schrijf de aftrekking en het antwoord.
\end{exercise}

\begin{exercise}[$\star$]\label{exo:g2:subtraction:7}
Controleer met de tweelingoptelling: $85 - 39 = 46$; klopt dat?
$214 - 68 = 156$; klopt dat?
\end{exercise}

\begin{exercise}[$\star$]\label{exo:g2:subtraction:8}
Kies bij elk verhaaltje $+$ of $-$: ``$47$ vogels, $19$ vliegen weg'';
``$47$ appels, er worden $19$ bij geplukt''; ``Sam heeft $47$, Lea heeft
$19$: hoeveel heeft Sam er meer?''.
\end{exercise}

\begin{exercise}[$\star$]\label{exo:g2:subtraction:9}
Vul het gat in: $63 - \;?\; = 45$. (Tel door van $45$ naar $63$.)
\end{exercise}

\begin{exercise}[$\star\star$]\label{exo:g2:subtraction:10}
Een spelletje: denk aan een getal, tel er $30$ bij op en haal er dan
$18$ af. De uitkomst is $58$. Aan welk getal dacht je? (Reken terug!)
\end{exercise}
